% ML4H 2026, proceedings track (8 pages). Revised for the author response.
\documentclass[pmlr,twocolumn,10pt]{jmlr}

\mlhtrack{proceedings}

\newif\iffinal
\finalfalse  % \finaltrue for camera-ready

\iffinal
    \ifmlhneedspmlr
      \jmlrvolume{XXX}
      \jmlryear{2026}
    \fi
    \ifmlhfindings \jmlrproceedings{}{ML4H 2026 - Findings Track}\fi
    \ifmlhdemo     \jmlrproceedings{}{ML4H 2026 - Demo Track}\fi
    \jmlrworkshop{Machine Learning for Health (ML4H) 2026}
\else
    \jmlrproceedings{}{Submitted to ML4H 2026: \mlhtrackname}
    \jmlrworkshop{Machine Learning for Health (ML4H) 2026}
\fi

\urlstyle{same}

\usepackage{booktabs}
\emergencystretch=3em
\usepackage{siunitx}
\usepackage[switch]{lineno}

\theorembodyfont{\upshape}
\theoremheaderfont{\scshape}
\theorempostheader{:}
\theoremsep{\newline}
\newtheorem*{note}{Note}

\title[Auditing the Auditor]{Auditing the Auditor: Four Failure Modes of
Confound Detection in Clinical Machine Learning}

\author{Anonymous Author(s)}

\linenumbers

\begin{document}

\maketitle

\ifmlhdemo\else
\begin{abstract}
Clinical machine learning results are routinely undermined by confounds that
leave no trace in the reported metrics. We built five automated detectors for
such confounds and then audited them the way they audit a pipeline. The audit
yields an initial catalogue of four ways a diagnostic can fail whether or not
its target confound is real, organized by the stage of a diagnostic at which
they arise: the statistic, the decision rule, the controls, and the design
process. Types 1 to 3 have measured instances in our own detectors, Type 1 also
has a structural instance in third-party evaluation code, and Type 4 is
supported only counterfactually. We do not claim the catalogue is complete.

The measured instances are mostly unflattering to our detectors. A diagnostic
with no undecidable state reported false cleanliness five times in our own code.
A conventional agreement threshold sat on the population value it had to
separate from ($\kappa=0.602$ against $0.60$), so a 500-patient audit is close
to a coin flip. At full scale our pre-registered corpus-size control recovers
55\% to 69\% of the pretraining-leakage detector's effect, which makes that
effect confounded with corpus size. The static detector for out-of-distribution
selection recognised none of 727 third-party files and returned
\textsc{indeterminate} on the files of all 53 sites we labelled by hand. A
falsy-as-clean pattern was reachable and unguarded in 4 of 48 applicable
sampled sites across 25 repositories (exact 95\% interval 0.02 to 0.20), or in
16 of 48 (0.20 to 0.48) if NaNs that arise only from diverged training count,
with no evidence that it concentrates in code that renders verdicts. This is a
critique, with negative results, of verification tooling.
\end{abstract}
\begin{keywords}
confound detection, reproducibility, clinical machine learning, label shift,
data leakage, model selection, negative results
\end{keywords}
\fi

\ifmlhneedsstatements
\paragraph*{Data and Code Availability}
We use three credentialed critical-care datasets from
PhysioNet~\citep{goldberger2000physiobank}: the PhysioNet 2019 Sepsis Challenge
corpus~\citep{reyna2020early}, MIMIC-IV~v3.1~\citep{johnson2023mimiciv} and
eICU-CRD~v2.0~\citep{pollard2018eicu}, each available to researchers who
complete the required training and sign the data use agreement. Our code is
available in an anonymized repository, \url{https://anonymous.4open.science/r/reboot-8568/}. It
contains all detector code and tests, the pre-registration and outcome
documents, run logs, and every aggregate result (confusion matrices, $\kappa$
values, confidence intervals, $t$-statistics); author-identifying commit history
is removed for review. Per-stay derived caches are not released because they are
individual-encounter artifacts of restricted data; the repository documents how
to regenerate them under the reader's own credentials, and the reproduction path
does not depend on them.

\paragraph*{Institutional Review Board (IRB)}
This work uses only de-identified, publicly available datasets obtained under
their respective data use agreements, and involves no interaction with human
subjects and no access to identifiable private information. It therefore does
not constitute human subjects research and did not require IRB approval.
% TODO: confirm IRB wording with institution.
\fi

\section{Introduction}
\label{sec:intro}

A clinical machine learning result can be wrong in a way that nothing in the
paper reveals. The model is sound, the code runs, the numbers are real, and the
finding is an artifact of how the data was assembled. Systematic reviews of
COVID-19 prediction models found essentially none fit for clinical
use~\citep{wynants2020prediction,roberts2021common}, and the failures were
overwhelmingly of this kind. Leakage in particular has been argued to underlie a
reproducibility crisis across fields that use machine
learning~\citep{kapoor2023leakage}.

The natural response is tooling: if a confound leaves a detectable signature,
write a detector and run it before publishing. We did this for five confound
types we had personally encountered: (1) label-definition shift between sites,
(2) pretraining leakage into a ``zero-shot'' target, (3) a hyperparameter swept
on out-of-distribution data, (4) a physiological constraint imposed on a
variable that was computed from it, and (5) a cross-site score difference
driven by which variables were recorded. Each detector was checked against a
positive case and a false-positive control, with ground truth taken from real
historical defects in our own repository wherever possible.

Taking the detectors seriously meant auditing them the way they audit a
pipeline, and the audit found more than the detectors did. A diagnostic can
fail independently of whether its target confound exists, in ways that are
invisible in the one thing a detector reports, a clean confusion matrix. This
paper is the audit. Its contributions are:
(i) an initial catalogue of four failure modes organized by the stage at which
they arise (\sectionref{sec:taxonomy});
(ii) the evidence for each type, with its strength stated per type
(\tableref{tab:evidence}) and quantified with exact intervals;
(iii) additional pre-registered experiments reported as found, including those
that go against our own detectors; and
(iv) a minimal reporting practice for diagnostics. The paper is a critique and
a set of negative results about verification tooling, and it does not claim the
catalogue is complete or that the detectors are established.

\section{Related Work}
\label{sec:related}

\citet{damour2022underspecification} show that a training pipeline can be
underspecified: many models fit equally well while behaving differently under
shift. We ask whether the \emph{diagnostic} is underdetermined by its own
validation. A detector reporting TN$=$2 does not tell you whether it could have
fired. The closest structural analogue is \citet{salaudeen2025misspecified},
who give conditions under which a domain generalization benchmark can detect
what it claims to and audit real benchmarks against them. A misspecified
benchmark yields an uninformative number; a vacuous diagnostic yields an
affirmatively wrong one, because ``did not detect'' reads as ``nothing to
detect''. The confounds we target are instances of shortcut
learning~\citep{geirhos2020shortcut} and hidden
stratification~\citep{oakdenrayner2020hidden}; checklists and
appraisals~\citep{kapoor2023leakage,wynants2020prediction,roberts2021common}
catalogue them for human readers, whereas a detector is executable and can be
wrong in a way one can measure.

Static detection of leakage in data science code is the closest tooling.
\citet{yang2022leakage} detect overlap, pre-processing and multi-test leakage in
notebooks, and \citet{subotic2022leakage} give a lightweight analysis for
development time. Multi-test leakage, reusing test data for model selection, is
near detector~3's confound. Their tools target notebook patterns around
train/test splits; detector~3 targets site-level sweep functions. We ran the
former on our own fixtures (\sectionref{sec:generalize}). AIF360~\citep{bellamy2019aif360}
emits verdicts rather than numbers for a human to read, and we read it for our
own failure modes.

\section{The Detectors and the Validation Protocol}
\label{sec:methods}

\paragraph{Detectors.} Detector~1 computes Cohen's $\kappa$~\citep{cohen1960kappa}
on a fixed 500-patient audit subset between two operationalizations of the
outcome and flags below the conventional Landis--Koch
boundary of $0.60$~\citep{landis1977measurement}; the reference definition is
window-mode SOFA per Sepsis-3~\citep{singer2016sepsis3}. Detector~2 trains a
small masked-reconstruction probe on a source site with a controlled fraction of
target stays leaked in and tests whether probe loss on held-out target data
falls, with a one-sided paired $t$ within seed ($df=4$ for 5 seeds). Detector~3
is static analysis: via Python's abstract syntax tree it traces whether a sweep
function's evaluation call receives a loader derived from the held-out
distribution. Detector~4 traces variable lineage through the same machinery,
with guard-condition reasoning, to find constraints imposed on variables derived
from those constraints. Detector~5 compares per-component availability ratios
and a composition-controlled aggregate across sites, and flags when the
availability ratio exceeds $2.0$ and the composition gap ratio exceeds $0.30$.

\paragraph{Ground truth.} The positive case is a real historical defect where
possible: detector~3's is a buggy sweep from our git history, kept with its
corrected version; detector~4's is a genuine Severinghaus circularity in our
PhysioNet loader; detector~1's is ICD-coded sepsis against SOFA-window scoring.
Only detector~2's leakage and detector~5's ablation are injected, because the
quantity detected is a dose.

\paragraph{Controls and margins.} Every detector ships a false-positive
control. For a stochastic detector, a control's \emph{margin} is
$m=|\tau-s|/\sigma$, the distance of its statistic $s$ from the threshold $\tau$
in standard deviations $\sigma$ of the statistic at the operating audit size. By a
rule fixed before we computed it, a control is \emph{discriminating} if
$m\le 3$ and a \emph{ceiling} control otherwise: one that passes because it is
far from the boundary and cannot fail. Every stochastic detector reports an
uncertainty measure or renders \texttt{MISSING}; detectors 3 and 4 are
deterministic and are excluded explicitly.

\paragraph{Pre-registration and scale.} Detector~5's variant rule was committed
29 minutes before the commit adding its external-validation code
(\appendixref{apd:prereg}). The additional experiments reported here were
pre-registered in committed documents before they ran, in order, with the
commit hashes in \appendixref{apd:prereg}. Every result states its scale. The
\emph{internal} suite runs at demonstration scale: detector~1 on the eICU demo
cohort, detector~2 on PhysioNet Site~A${\to}$Site~B at 900 nominal stays per
site (602 to 651 source stays remain after the loader's filters), detector~5 on
1200 stays per arm. \emph{External} runs use full MIMIC-IV~v3.1 and
eICU-CRD~v2.0 and are marked full-scale.

\section{Internal Validation}
\label{sec:results}

\begin{table*}[t]
\floatconts
  {tab:main}
  {\caption{\textbf{Internal} validation suite at demonstration scale, with exact
   (Clopper--Pearson) 95\% intervals, and control margins. Controls are counted
   as discriminating / ceiling / margin unknown over both scales. The pooled row
   is reference only: the cases are heterogeneous and partly correlated, and we
   do not read it as performance. Detector~5's row is a false negative.}}
  {\footnotesize\setlength{\tabcolsep}{3pt}%
  \begin{tabular}{@{}lccccp{3.3cm}@{}}
  \toprule
  \bfseries Detector & \bfseries TP/FP/FN/TN & \bfseries Recall [95\% CI] &
    \bfseries FPR [95\% CI] & \bfseries Controls & \bfseries Uncertainty, external\\
  \midrule
  1 label shift & 1/0/0/1 & 1/1 [0.025, 1.000] & 0/1 [0.000, 0.975] & 1/8/3 & full MIMIC-IV and eICU-CRD\\
  2 leakage     & 3/0/0/1 & 3/3 [0.292, 1.000] & 0/1 [0.000, 0.975] & 0/0/1 & paired $t$, $df=4$; full MIMIC-IV$\to$eICU\\
  3 OOD select. & 1/0/0/1 & 1/1 [0.025, 1.000] & 0/1 [0.000, 0.975] & n/a & 0 of 727 files recognised\\
  4 circularity & 1/0/0/8 & 1/1 [0.025, 1.000] & 0/8 [0.000, 0.369] & n/a & no external positive\\
  5 missingness & 0/0/1/1 & 0/1 [0.000, 0.975] & 0/1 [0.000, 0.975] & 0/0/2 & external ablation\\
  \midrule
  Pooled (ref.\ only) & 6/0/1/12 & 6/7 [0.421, 0.996] & 0/12 [0.000, 0.265] & 1/8/6 & precision 6/6 [0.541, 1.000]\\
  \bottomrule
  \end{tabular}}
\end{table*}

\tableref{tab:main} gives the internal suite. Its totals are $\mathrm{TP}=6$,
$\mathrm{FP}=0$, $\mathrm{FN}=1$, $\mathrm{TN}=12$, but the denominators are
small (6, 7 and 12), the intervals are wide, and the cases are heterogeneous, so
the totals say little about how any detector would perform elsewhere. Detector~1's
$\mathrm{TN}=1$ is the corrected control, SOFA window-mode against single-point
scoring, a ceiling control by our rule ($m=3.9$); the first version, which scored
an array against itself, was replaced before the matrix was computed
(\sectionref{sec:measured}). Detector~4's eight true negatives are non-circular
constraint/loader pairs across the three databases. The external controls are
not folded into the matrix: five of the six full-scale detector~1 controls are
ceiling controls, so adding them would raise TN on tests that discriminate
nothing, the denominator inflation this paper argues against.


\begin{table*}[t]
\floatconts
  {tab:depth}
  {\caption{Validation depth per detector. Controls are counted over both scales;
   ``disc./ceil./unk.'' classifies them by the margin rule of
   \sectionref{sec:methods}. Detectors 3 and 4 are deterministic.}}
  {\footnotesize\setlength{\tabcolsep}{3pt}%
  \begin{tabular}{@{}lcccp{2.5cm}p{2.1cm}p{2.0cm}@{}}
  \toprule
  \bfseries Det. & \bfseries Positives int./ext. & \bfseries Controls int./ext. &
    \bfseries Disc./ceil./unk. & \bfseries External ground truth &
    \bfseries Uncertainty & \bfseries Baseline, manual labels\\
  \midrule
  1 & 1 / 2 & 6 / 6 & 1 / 8 / 3 & two Sepsis-3 operationalizations & 5000-resample bootstrap & prevalence ratio (inverted); n/a\\
  2 & 3 / 3 & 1 / 1 & 0 / 0 / 1 & controlled injection (dose) & paired $t$, $df=4$; null 1/20 & three standard checks; n/a\\
  3 & 1 / 0 & 1 / 0 & n/a & historical defect; 727 files, 53 sites & deterministic & Yang et al., zero; 53 sites\\
  4 & 1 / 0 & 8 / 0 & n/a & no external positive found & deterministic & none found; n/a\\
  5 & 1 / 3 & 1 / 2 & 0 / 0 / 2 & controlled ablation (dose) & seed flag rate & out of scope; n/a\\
  \bottomrule
  \end{tabular}}
\end{table*}

\section{A Catalogue of Diagnostic Failure Modes}
\label{sec:taxonomy}

We define a \emph{diagnostic} as a procedure that audits an ML pipeline for one
confound and returns a verdict. Our catalogue is initial: it organizes the
failures we found by the stage of a diagnostic at which they arise, and we claim
completeness only with respect to that stage model. A diagnostic has four
stages: (i) a statistic $S$ computed from the audited artifacts; (ii) a decision
rule mapping $S$ to a verdict; (iii) controls that validate the pair; and (iv)
the design process that produced them. The types are:

\textbf{Type 1, vacuous pass} (stage ii--iii interface). A check that never
reached a verdict is scored as ``clean''. \emph{1a, falsy-as-clean}: the verdict
function is total on a domain where $S$ is partial, so an undefined result
(\texttt{nan}, \texttt{None}, an empty container) is consumed by a boolean test.
\emph{1b, degenerate control}: a control with margin $m>3$ (or an identity) is
counted as a true negative though it cannot fail.

\textbf{Type 2, structural non-discriminability} (stage ii). For two
implementations that a domain convention deems interchangeable, the population
value $\theta^*$ of the statistic satisfies $|\tau-\theta^*|/\sigma_{\text{audit}}\le 1$.
The flag probability, approximately $\Phi((\tau-\theta^*)/\sigma_{\text{audit}})$, is then
far from 0 and 1 at that audit size. A larger audit would make the verdict
deterministic but still decided by a margin far smaller than the spread of
$\theta^*$ across legitimate implementations; the issue is the threshold's
position relative to that band, not sample size.

\textbf{Type 3, granularity or pairing mismatch} (stage i). The statistic's
aggregation level or pairing differs from the structure of the confound.

\textbf{Type 4, post-hoc fitting risk} (stage iv). A rule changed after
observing its failure gives contaminated evidence unless it was fixed in advance.
Type~4 is a process hazard; we did not observe it as a failure of a shipped
detector.

The catalogue does not cover: an invalid reference or null distribution,
non-independent audit draws, plain implementation bugs, or a miscalibrated
threshold that is not of Type~2. All four types are compatible with a perfect
confusion matrix, which is why we report the matrix and the catalogue together.
\tableref{tab:evidence} states how much each type is supported.

\begin{table}[htbp]
\floatconts
  {tab:evidence}
  {\caption{Evidence strength per type. ``Quantified'' names the measurement.}}
  {\footnotesize\setlength{\tabcolsep}{2pt}%
  \begin{tabular}{@{}lp{2.0cm}p{1.9cm}p{2.3cm}@{}}
  \toprule
  \bfseries Type & \bfseries Internal & \bfseries External & \bfseries Quantified\\
  \midrule
  1 & 5 in our code & structural, 4/48 to 16/48 sites & exact CIs; margins\\
  2 & detector~1; detector~5 (exploratory) & none & bootstrap; consistency check\\
  3 & 2 initial false negatives; 1 redundant gate & none & seed flag rates\\
  4 & counterfactual only & none & pre-registration record\\
  \bottomrule
  \end{tabular}}
\end{table}

\section{The Failure Modes, Measured}
\label{sec:measured}

\subsection{Type 1: five vacuous passes in our own code}

When a check cannot reach a judgement, because the metric is undefined, the
pattern is absent or the control is too easy to fail, that outcome silently
becomes ``no confound found'' and is counted as evidence of specificity. We hit
this five times, independently:

\begin{enumerate}
\itemsep0em
  \item \emph{(1b)} Detector~1's negative control passed the same audit array as
    both arguments to $\kappa$, giving $1.000$ by construction.
  \item \emph{(1b)} Detector~1's suspicion-window width controls sit at $\kappa$
    $0.93$ to $1.00$; reported bare, the external $\mathrm{TN}=3$ implies three
    tests when only one discriminates.
  \item \emph{(1a)} Detector~3's \textsc{indeterminate} verdict on third-party
    repositories, folded into TN, would have produced $\mathrm{TN}=3$ from files it
    never analysed.
  \item \emph{(1a)} A baseline's \texttt{train\_test\_gap}: one seed's in-domain
    AUROC was undefined, \texttt{nan > threshold} was \texttt{False}, and the
    abstention was scored as a correct silence.
  \item \emph{(1b)} Detector~2's 0\% control compares an arm with itself, so its
    relative delta is exactly zero and $\mathrm{TN}=1$ holds by construction.
\end{enumerate}

The fix is mechanical: every check returns a decidability flag with its verdict,
and undecidable cases are reported in their own column. A sixth, different
instance was stale evidence: two committed logs still reported a retracted
composition figure and a passing verdict for the case this paper
reports as a false negative; they are regenerated and a test now scans for the
retracted wording.

\paragraph{Detector 2: corpus size, the null, and values versus missingness.}
Because the 0\% control holds corpus size fixed, it could not show what a
corpus-size control shows. We pre-registered one (arm A2: extra \emph{source}
stays in place of target stays). As a fraction of the shipped detector's effect,
its effect $r_2/r_1$ is $0.31$ at 5\% and $0.48$ at 20\% leakage on the PhysioNet
demonstration pair and \textbf{0.55 and 0.69} at full scale (MIMIC-IV$\to$eICU),
both full-scale values above the pre-registered line of $0.5$. By our own rule,
detector~2's effect is confounded with corpus size. Target exposure still adds
beyond volume (the shipped arm against A2: paired $t=9.0$ and $15.4$ at full
scale), and the detector is unchanged; what it establishes is narrower than we
claimed. We also ran an exploratory null of 20 independent paired five-seed
comparisons of two no-leak arms: 1 of 20 flagged (exact 95\% interval 0.001 to
0.249), consistent with the nominal 5\% but with an upper bound of 25\% and a
harder null than the detector's own shared-seed comparison
(\appendixref{apd:expl}). On the demonstration pair, the pre-registered arms
separating target values from target observation pattern gave a missingness
share $R=0.02$ and $-0.11$ at 5\% and 20\% (``values-dominated'' by the rule);
our pre-stated prediction, mixed with $0.25<R<0.6$, was wrong. The mask-transplant
arm is under-filled (0.215 of cells observed against 0.301 in the target), so a
missingness contribution is not excluded, and at full scale $R$ is undetermined
because the available caches lack the raw pre-forward-fill series.

\subsection{Type 2: a threshold on the value it must separate}

On full MIMIC-IV ($74{,}829$ stays) the whole-cohort $\kappa$ between SOFA
window-mode and single-point scoring, two defensible Sepsis-3 operationalizations,
is $\mathbf{0.602}$ (95\% CI $[0.597, 0.607]$), $0.002$ above the $0.60$ boundary
and $0.06$ standard deviations of the audit-draw distribution from it. A
500-patient audit therefore flags $48.4\%$ of the time; a normal approximation
with the draws' own spread gives $0.475$, a consistency check rather than a
prediction, since both use the same draws.

Two estimands must not be compared. On the 117-stay demonstration cohort the
audit subset is the whole cohort, so the audit-draw flag probability is
degenerate ($0.000$); the cohort-bootstrap probability is $0.216$ there and
$0.241$ at full scale. Like for like it barely moved with $640\times$ more data,
and the demonstration estimate ($0.651$, CI containing $0.602$) could not
resolve the near-miss. The full-scale value does: the population sits on the
threshold. The same implementation pair gives $\kappa=0.826$ on eICU-CRD, a
ceiling control ($m=5.39$), so legitimate variation across databases (0.224)
dwarfs the 0.002 by which the MIMIC control misses. We did not move the threshold
(that would be Type~4). The value $0.60$ is a convention with no clinical
derivation; we offer as a prescription for new detectors, fixed before data,
that a negative control be a pair a domain expert accepts as interchangeable and
that a threshold sit clear of the legitimate-variation band of such pairs. The
full-scale eICU positive sits at $m=2.59$ and the MIMIC positive at $12.1$.

An exploratory second instance: detector~5's composition ratio on PhysioNet
Site~A vs.~B (gate $0.30$) has mean $0.283$ (SD $0.082$) over five seeds at
$n=1200$ stays per arm and flagged 2/5, mean $0.272$ (SD $0.027$) and 0/5 at
$n=4000$, and 3/5 at $n=900$ and $1500$. A normal approximation gives flag
probabilities $0.42$ and $0.15$ at the first two sizes; this is consistent with
the counts, with five seeds per cell and the same seeds supplying both
estimate and count.

\subsection{Type 3: granularity and pairing}

Two of the five detectors first produced false negatives for the same reason:
detector~5 aggregated at the wrong granularity and detector~2 compared unpaired
where the comparison is paired, so its unpaired null spread is comparable to
the effect and it is reported as a paired $t$. At full scale detector~5's
composition gate binds nowhere, so the conjunction collapses into the availability
condition alone and one of two gates is decorative. Restricting detector~5 to the
components present in every database (MAP and Henderson--Hasselbalch) flags the
PhysioNet flagship in 5 of 5 seeds, but the ratio is unstable on the same-site
control (mean $2.9$, range $0.16$ to $12.8$ at $n=900$), so the gate is
decorative on PhysioNet too. The composition quantity is a ratio of two gaps,
not a share.

\subsection{Type 4: what pre-registration bought}

We did not observe Type~4 as a failure of a shipped detector; the evidence is
counterfactual. We pre-registered a rule for choosing between two variants of
detector~5 before running either on external data. The alternative had higher
demonstration-scale recall ($0.73$ against $0.60$) and was tempting; the
pre-registered criterion, a strictly lower detection floor, rejected it. The
pre-registration also excluded in advance one case with unknown ground truth,
the only one on which the variants diverge: scored, it alone would have decided
the comparison. Fitting detector~3's recognition vocabulary to the corpus (19 to
60 names) was done as a labelled arm and did not help.

\section{Generality and External Evidence}
\label{sec:generalize}

\subsection{Detector 3 outside its home idiom}

Over 18 pinned public repositories, each associated with a paper (735 Python
files, 727 parsed, 8 unparseable and reported as such), detector~3 recognised a
sweep in \textbf{none}, with its shipped vocabulary and with a vocabulary fitted
to those repositories. To measure sensitivity with a denominator we enumerated
candidate selection sites with patterns frozen before the scan (103 sites),
sampled 53, and labelled each with a rubric committed before any site was read
and before the detector ran on it. None was contaminated; one was a disclosed
oracle (a maximum over per-epoch test accuracy, printed as ``Target
Validation''), one was ambiguous, 21 were sound and 30 were not selections.
Sensitivity is therefore undefined, and undisclosed contamination was absent
from 0 of 23 real selection sites (exact 95\% interval 0.000 to 0.148; the sites
lie in 15 files). The detector returned \textsc{indeterminate} on the files of
all 53, under both vocabularies, so it also issues no clean verdict. Detector~3
is a guard for one's own sweeps, not a retrospective auditor. The scan covers
Python files only: 126 notebooks were not scanned.

We built the tool of \citet{yang2022leakage} in Docker, after three changes to
its build environment, and ran it on detector~3's buggy and corrected fixtures.
It reports zero in all three categories on both and recognises no
train/validation/test model pairing in either file. Our fixtures are not in the
format the tool targets, so this shows that it does not apply to our idiom after
those changes, not that it fails in general.

\subsection{Type 1 beyond our code}

We read 25 repositories (the 18 above, AIF360~\citep{bellamy2019aif360} and six
new ones chosen by a star-ranked rule fixed before any query) for the
falsy-as-clean pattern. Patterns for four kinds of site were frozen first (236
sites, 88 sampled) and each sampled site was labelled GUARDED, UNREACHABLE,
REACHABLE\_UNGUARDED or not applicable. The convention for NaNs that arise only
when training diverges was \textbf{not fixed in advance}, so we report both
readings: reachable and unguarded in $4/48$ applicable sites (exact 95\%
interval $[0.023, 0.200]$) if such NaNs are outside supported use, and in
$16/48$ ($[0.204, 0.484]$) if they count. The four positives under the narrow
reading are in three files (\appendixref{apd:positives}). Splitting applicable
sites into verdict or selection code and other code gives $1/27$ against $3/21$
(Fisher exact $p=0.31$) under the narrow reading and $11/27$ against $5/21$
($p=0.35$) under the broad one. The direction reverses with the convention, so
we found \emph{no evidence that the pattern concentrates in code that renders
verdicts}; it appears in selection, filtering and coercion code alike. The
external evidence supports Type~1 only.

DomainBed~\citep{gulrajani2021search} shows the structural point. Its query
library maps \texttt{nan} to $-\infty$ inside its sorted routine, yet its
\texttt{argmax}, used to pick the winning run, is a bare \texttt{max}: a
\texttt{nan} score wins or loses by list position. Its leave-one-out selection
guards missing environments with \texttt{any([v == -1 for v in val\_accs])},
vacuously false on an empty list. Neither is reachable with the shipped datasets,
which have at least three environments, so we report them as structurally
present and not currently exploitable.

\section{Comparison with Standard Practice}
\label{sec:baselines}

We ran three standard checks (external floor, $k$-fold instability, train/test
gap) on the same trained models as detector~2's leakage injection. At
demonstration scale they flag at 0\% leakage in 4/5, 5/5 and 4/4 decidable seeds.
At full scale, $k$-fold instability is silent at every level, the floor is silent
except 2/5 seeds at 100\%, and only the gap fires, at 0\% too. None carries
information about this scenario, in different ways at the two scales
(\appendixref{apd:base}). We do not claim that detector~2 beats them: its
effect is partly corpus size.

Downstream change is small. Cross-site AUROC changes by $+0.0007$
$[-0.0010, 0.0025]$, $+0.0027$ $[-0.0009, 0.0062]$ and $+0.0037$
$[0.00003, 0.0074]$ at 5\%, 20\% and 100\% leakage on the full held-out
demonstration target (13{,}879 to 13{,}888 stays, 285 to 289 positives), and by $+0.0007$,
$0.0000$ and $-0.0025$, all intervals including zero, at full scale (about 2{,}900
positives), against a 0\% source-to-target gap of $0.059$ to $0.142$. The
exposure detector~2 measures does not produce a measurable downstream change in
our setting, so its practical importance is not shown. A prevalence-ratio check
for detector~1 is inverted: it ranks legitimate cross-site comparisons above
real label mismatches (AUROC for picking mismatches $0.059$ and $0.125$ at the
two scales). Detector~5's baseline would be the naive gap it decomposes, so it
is not independent and was not run; we found no baseline for detector~4 in two
searches.

\section{Discussion}
\label{sec:discussion}

A clean confusion matrix is not evidence that a detector works: every failure
above is compatible with \tableref{tab:main}. The minimum a diagnostic should
report with its verdict is an explicit undecidable count, the margin of each
control, an uncertainty measure at the operating point, and the decision rule's
provenance. The recommendation that stands from the external read is to give
every diagnostic an explicit undecidable state. Scale changed what two
diagnostics do (\sectionref{sec:baselines}, detector~5's gates), so every figure
here carries its scale.

\section{Limitations}
\label{sec:limitations}

The same group built and audited every detector; the external scans, the
hand labelling and the owner's spot-check narrow this circularity but do not
remove it, and the spot-check is an agreement figure, not a blind relabel.
Types 2 to 4 have no external instance, Type~2 rests on one detector plus an
exploratory second, and Type~3 on two early false negatives and one redundant
gate, so it is the type most likely to be revised. The internal suite is small
and heterogeneous (\tableref{tab:main}). Detector~3 has no demonstrated
sensitivity outside its home idiom, and its scan and the Type~1 read cover
Python files only. Detector~2's effect is confounded with corpus size at full
scale; the missingness arms are demonstration-scale only; and detector~5's
external component set is analogous, not identical. Baselines cover detector~2
and a few others only. Full-scale detector~2 and baseline runs use caches of
74{,}607 and 130{,}446 stays, slightly different from the 74{,}829 and 132{,}900
used for detector~1. The formal definitions here were written by us and have not
been reviewed by outside domain experts, a step we would welcome. We did not
change any detector or threshold after observing its failure.

\section{Conclusion}
\label{sec:conclusion}

We built five confound detectors and audited them. The audit gives an initial
catalogue of four failure modes, organized by the stage at which they arise,
with evidence of different strength per type (\tableref{tab:evidence}), most of
it unfavourable to our own detectors. The practical recommendation is narrow and
cheap: give every diagnostic an explicit undecidable state, check that every
control has a margin and is not an identity, report uncertainty at the operating
point, and fix the decision rule before looking.

\acks{Acknowledgments go here \emph{but should only appear in the camera-ready
version of the paper if it is accepted}.}

\bibliography{ref}

\appendix

\section{Reproduction}\label{apd:repro}

The released repository reproduces \tableref{tab:main} from a clean clone with
no clinical data present: detectors 3 and 4 run live as static analysis, and
detectors 1, 2 and 5 load committed aggregate results. The accompanying test
suite runs 116 tests, of which 7 \emph{skip} with the explicit reason
\texttt{PhysioNet data not present} rather than passing: a test that cannot run
has not succeeded. Per-stay arrays derived from MIMIC-IV and eICU-CRD are not
released.

\section{Pre-registration}\label{apd:prereg}

The pre-registration and outcome documents are in the anonymized repository,
with anonymized patches of the two commits that fix the order for detector~5
(\texttt{detectors/preregistration\_evidence/}): the pre-registration
(\texttt{63f5d95a8a8e}, 2026-08-06 00:34) and its direct child adding the
variant rules and external-validation code (\texttt{1600d91071cf}, 01:03).
The experiments in this revision were pre-registered in the repository before
they ran: the main document (\texttt{25fc617}), the detector~3 and Type~1 site
patterns and rubrics (\texttt{b845cc0}, \texttt{4c461eb}, \texttt{bc365cf},
\texttt{073d042}, \texttt{879d885}), addendum 2 (\texttt{d14687a}) and addenda 3
and 3b for the full-scale tier (\texttt{4918770}, \texttt{69148a9}). Dates are
author-supplied metadata; the full hashes are the check.

\section{Exploratory analyses}\label{apd:expl}

\emph{Null for detector~2.} A replicate is an independent paired five-seed
comparison of two no-leak arms, each with a fresh data draw and model seed,
flagged by the shipped rule; 1 of 20 flagged (exact 95\% interval 0.001 to
0.249). The arms use independent model seeds, harder than the detector's own
shared-seed comparison; 3 of 20 replicates had $|t|>2.132$ (probability 0.075 of
three or more at a 5\% level). \emph{Full held-out target (5a-prime).} Models
retrained under the identical seeds; probe loss agrees with the earlier run to
$3\times10^{-8}$. Target stays 13{,}879 to 13{,}888 per seed, 285 to 289
positives; the 0\% in-domain minus target gap is $0.142$ $[0.029, 0.254]$.
\emph{Detector~5.} Five-seed external run, variant A: precision $9/9$, FPR
$0/10$, flag rates $0.2$, $0.8$, $0.8$ at 50\%, 80\%, 95\% ablation. All items in
this appendix are exploratory.

\section{Baselines at the two scales}\label{apd:base}

\begin{center}\footnotesize
\begin{tabular}{@{}lcccc@{}}
\toprule
\bfseries Seeds flagging & \bfseries 0\% & \bfseries 5\% & \bfseries 20\% & \bfseries 100\%\\
\midrule
\multicolumn{5}{@{}l}{\emph{Demonstration scale (flags/decidable, 5 seeds)}}\\
$k$-fold instability & 5/5 & 5/5 & 5/5 & 5/5\\
train/test gap & 4/4 & 4/4 & 4/4 & 4/4\\
external floor & 4/5 & 4/5 & 4/5 & 3/5\\
detector 2 (paired $t$) & --- & flag & flag & flag\\
\midrule
\multicolumn{5}{@{}l}{\emph{Full scale (flags/5 seeds, Sepsis-3 label)}}\\
$k$-fold instability & 0/5 & 0/5 & 0/5 & 0/5\\
train/test gap & 5/5 & 4/5 & 5/5 & 4/5\\
external floor & 0/5 & 0/5 & 0/5 & 2/5\\
detector 2 (paired $t$) & --- & $-3.7$ & $-15.3$ & $-35.1$\\
\bottomrule
\end{tabular}\end{center}

\section{Type 1 positives}\label{apd:positives}

Reachable and unguarded under the narrow reading (repository, file, line, pinned
commit): tableshift, \texttt{tableshift/\allowbreak{}datasets/\allowbreak{}college\_scorecard.py}:170,
\texttt{fca9429814} (a NaN target becomes 0 before NaN rows are dropped);
fairseq2, \texttt{src/\allowbreak{}fairseq2/\allowbreak{}data/\allowbreak{}parquet/\allowbreak{}arrow\_transform.py}:288 and :306,
\texttt{7f06d6f4f5} (null rows pass a length filter); robustdg,
\texttt{utils/\allowbreak{}scripts/\allowbreak{}utils.py}:579, \texttt{3eee1730ae} (\texttt{max} of two
losses returns the finite one when the other is NaN, defeating the
\texttt{isfinite} check on the next line). Reachability for the first is inferred
from the dataset's missing-value settings and was not executed.

\end{document}
